%! Setting Up a Test for the Difference of Two Population Means — Unit 7, Lesson 7.8 — Homework
\documentclass[10pt]{article}
\usepackage{apstats-article}
\usepackage{apstats-boxes}
\usepackage{grading}
\gradingassignment{apstats-u07-l08-hw}
\gradingcourse{advanced_statistics}
\newcommand{\TallMath}[1]{$\displaystyle #1\rule[-1.4em]{0pt}{3.2em}$}

\begin{document}

\pageheader{Unit 7, Lesson 7.8}{Homework}

\vspace{0.12in}

\textit{For each problem: (a) name the inference method; (b) define $\mu_1$ and $\mu_2$ and
state $H_0$ and $H_a$; (c) verify conditions. Do \textbf{not} compute a test statistic.}

\begin{enumerate}

\item \textbf{Pain Relief.} \quad
A pharmaceutical company randomly assigns 24 patients to a new pain-relief drug (Group 1)
and 26 patients to a placebo (Group 2). Researchers believe the new drug produces a
\emph{lower} mean pain score on a 10-point scale. Both groups' data appear approximately
symmetric with no outliers.

\begin{enumerate}[label=\alph*.]
  \item Method: \makebox[5.5cm]{\smash{\gradedblank[id=hw1a, type=short, points=1, answer={{two-sample t-test for mu1 - mu2}}, check=text, display={\boldmath two-sample $t$-test for $\mu_1 - \mu_2$}]{5.3cm}{0.19in}}}
  \item $\mu_1 = $\par
        \gradedhide{\gradedhere{\gradedblank[id=hw1b1, type=rubric, points=1, display={\bfseries\boldmath mean pain score, new drug}]{\linewidth+6pt}{0.19in}}\writelines{9}}\par
        $\mu_2 = $\par
        \gradedhide{\gradedhere{\gradedblank[id=hw1b2, type=rubric, points=1, display={\bfseries\boldmath mean pain score, placebo}]{\linewidth+6pt}{0.19in}}\writelines{9}}

        $H_0$: \makebox[4.5cm]{\smash{\gradedblank[id=hw1b3, type=short, points=1, answer={{mu1 - mu2 = 0}}, check=text, display={\boldmath$\mu_1 - \mu_2 = 0$}]{4.3cm}{0.19in}}} \quad $H_a$: \makebox[4.5cm]{\smash{\gradedblank[id=hw1b4, type=short, points=1, answer={{mu1 - mu2 < 0}}, check=text, display={\boldmath$\mu_1 - \mu_2 < 0$}]{4.3cm}{0.19in}}}

        Justify the direction of $H_a$: \hfill\makebox[4.6in][l]{\gradedhere{\gradedblank[id=hw1b5, type=rubric, points=1, display={\bfseries\boldmath researchers expect the drug to lower mean pain: $\mu_1 < \mu_2$}]{4.6in}{0.19in}}}\\
                \gradedhide{\writelines{9}}
  \newpage
  \item Conditions:
        \begin{itemize}
          \boxguard
          \item Random: \hfill\makebox[4.6in][l]{\gradedhere{\gradedblank[id=hw1c1, type=rubric, points=1, display={\bfseries\boldmath patients randomly assigned to the groups --- met}]{4.6in}{0.19in}}}\\
                \gradedhide{\writelines{9}}
          \boxguard
          \item Normal ($n_1 = 24 < 30$, $n_2 = 26 < 30$): \hfill\makebox[3.8in][l]{\gradedhere{\gradedblank[id=hw1c2, type=rubric, points=1, display={\bfseries\boldmath both $n < 30$; both symmetric, no outliers --- met}]{3.8in}{0.19in}}}\\
                \gradedhide{\writelines{9}}
        \end{itemize}
\end{enumerate}
\vspace{0.08in}

\item \textbf{Wages.} \quad
An economist randomly surveys 50 workers in the technology sector (Group 1) and 45 workers
in the education sector (Group 2). She wants to determine whether the mean annual wages
\emph{differ} between the two sectors.

\begin{enumerate}[label=\alph*.]
  \item Method: \makebox[5.5cm]{\smash{\gradedblank[id=hw2a, type=short, points=1, answer={{two-sample t-test for mu1 - mu2}}, check=text, display={\boldmath two-sample $t$-test for $\mu_1 - \mu_2$}]{5.3cm}{0.19in}}}
  \item $H_0$: \makebox[4.5cm]{\smash{\gradedblank[id=hw2b1, type=short, points=1, answer={{mu1 - mu2 = 0}}, check=text, display={\boldmath$\mu_1 - \mu_2 = 0$}]{4.3cm}{0.19in}}} \quad $H_a$: \makebox[4.5cm]{\smash{\gradedblank[id=hw2b2, type=short, points=1, answer={{mu1 - mu2 != 0}}, check=text, display={\boldmath$\mu_1 - \mu_2 \ne 0$}]{4.3cm}{0.19in}}}
  \newpage
  \item Conditions:
        \begin{itemize}
          \boxguard
          \item Random: \hfill\makebox[4.6in][l]{\gradedhere{\gradedblank[id=hw2c1, type=rubric, points=1, display={\bfseries\boldmath random samples from both sectors --- met}]{4.6in}{0.19in}}}\\
                \gradedhide{\writelines{9}}
          \boxguard
          \item Normal: \hfill\makebox[4.6in][l]{\gradedhere{\gradedblank[id=hw2c2, type=rubric, points=1, display={\bfseries\boldmath $n_1 = 50,\ n_2 = 45 \ge 30$: CLT --- met}]{4.6in}{0.19in}}}\\
                \gradedhide{\writelines{9}}
        \end{itemize}
\end{enumerate}
\vspace{0.08in}

\item \textbf{Screen Time.} \quad
A pediatrician randomly selects 18 children in households with strict screen-time rules
(Group 1) and 21 children in households with no rules (Group 2). She believes children
with rules have a \emph{lower} mean daily screen time. A dotplot of Group 1 shows mild
left skew with no outliers; a dotplot of Group 2 shows approximately symmetric data.

\begin{enumerate}[label=\alph*.]
  \item Method: \makebox[5.5cm]{\smash{\gradedblank[id=hw3a, type=short, points=1, answer={{two-sample t-test for mu1 - mu2}}, check=text, display={\boldmath two-sample $t$-test for $\mu_1 - \mu_2$}]{5.3cm}{0.19in}}}
  \item $H_0$: \makebox[4.5cm]{\smash{\gradedblank[id=hw3b1, type=short, points=1, answer={{mu1 - mu2 = 0}}, check=text, display={\boldmath$\mu_1 - \mu_2 = 0$}]{4.3cm}{0.19in}}} \quad $H_a$: \makebox[4.5cm]{\smash{\gradedblank[id=hw3b2, type=short, points=1, answer={{mu1 - mu2 < 0}}, check=text, display={\boldmath$\mu_1 - \mu_2 < 0$}]{4.3cm}{0.19in}}}
  \item Conditions:
        \begin{itemize}
          \boxguard
          \item Random: \hfill\makebox[4.6in][l]{\gradedhere{\gradedblank[id=hw3c1, type=rubric, points=1, display={\bfseries\boldmath both groups randomly selected --- met}]{4.6in}{0.19in}}}\\
                \gradedhide{\writelines{9}}
          \newpage
          \boxguard
          \item Normal (both $n < 30$): \hfill\makebox[4.6in][l]{\gradedhere{\gradedblank[id=hw3c2, type=rubric, points=1, display={\bfseries\boldmath both $n < 30$; mild skew / symmetric, no outliers --- met}]{4.6in}{0.19in}}}\\
                \gradedhide{\writelines{9}}

                Can the test proceed? \makebox[1.5cm]{\smash{\gradedblank[id=hw3c3, type=short, points=1, answer=yes, check=text, display={Yes}]{1.3cm}{0.19in}}} \quad Explain:\par\vspace{5pt}
                \gradedhide{\gradedhere{\gradedblank[id=hw3c4, type=rubric, points=1, display={\bfseries\boldmath mild skew, no outliers: OK}]{\linewidth+6pt}{0.19in}}\writelines{9}}
        \end{itemize}
\end{enumerate}
\vspace{0.08in}

\newpage
\item \textbf{AP-Style.} \quad
\textit{A researcher randomly assigns 15 subjects to a cold-water immersion group and
15 subjects to a stretching group. He records recovery time (minutes) after exercise.
He tests $H_0: \mu_1 - \mu_2 = 0$ vs.\ $H_a: \mu_1 - \mu_2 \ne 0$. Both dotplots
show moderate skewness with one outlier each. Which statement about the Normal condition
is most accurate?}

\begin{enumerate}[label=(\Alph*)]
  \item Both conditions are met because $n_1 + n_2 = 30 \ge 30$.
  \item Both conditions are met; moderate skewness with one outlier per group is acceptable for $n < 30$.
  \item Both conditions are not met; outliers with $n < 30$ violate the Normal condition.
  \item The condition applies only to the group with the smaller sample size.
\end{enumerate}
\vspace{0.04in}
Answer: \makebox[1.5cm]{\smash{\gradedblank[id=hw4a, type=choice, points=1, answer=C]{0.4in}{0.2in}}}\rule[-0.22in]{0pt}{0.52in}

Brief justification:\par
\gradedhide{\gradedhere{\gradedblank[id=hw4b, type=rubric, points=1, display={\bfseries\boldmath both $n < 30$ with an outlier in each group}]{\linewidth+6pt}{0.19in}}\writelines{9}}

\end{enumerate}

\begin{extensionbox}
A student argues: ``Since the null hypothesis is always $\mu_1 - \mu_2 = 0$, a two-sample
$t$-test can never detect a specific numerical difference like $\mu_1 - \mu_2 = 5$.''
Is this correct? Explain when a researcher might hypothesize a non-zero difference and
how the null hypothesis would be written in that case.
\writelines{3}
\end{extensionbox}

\end{document}
